
\chapter{Literature Review}
\label{chap:literature_review}
In Uganda’s capital, Kampala, private shuttle and intercity transport services operate alongside other informal transport systems, serving passengers traveling between urban and upcountry destinations \cite{ndibatya2020kampala}, \cite{courtright2021uganda}. While these shuttle services are generally more comfortable and reliable than traditional public taxis, they still operate within a largely manual and semi-organized framework that presents challenges to both passengers and service providers \cite{howwemadeit2023ugandaTransport}, \cite{manwaring2021informal}, \cite{entebbe2026privateTransfers}.

In many cases, private shuttle operators rely on phone calls, referrals, and informal scheduling to manage bookings and coordinate trips \cite{howwemadeit2023ugandaTransport}. Departures are typically fixed within specific time windows, but seat allocation and passenger confirmations are often not managed through a centralized system. This can result in inefficiencies such as seat underutilization, customer no-shows, and difficulty in managing last-minute bookings \cite{manwaring2021informal}, \cite{masinde2024uganda}.

From the passenger perspective, accessing these services can also be challenging, as information about available routes, departure times, and seat availability is not always easily accessible in real time \cite{intalinc2021kampala}. In addition, courier services offered by shuttle operators are often managed manually, with limited tracking capabilities and reliance on direct communication between senders and operators \cite{masinde2024uganda}.

These limitations highlight gaps in the coordination and digital integration of private shuttle and courier services. Existing literature on transport systems suggests that mobile-based platforms with real-time booking, tracking, and communication features can significantly improve operational efficiency and service reliability, particularly in environments where transport systems remain semi-formal and manually coordinated. 

\section{Technologies Enabling Transport Digitization}
In recent years, the rise of mobile technology has revolutionized the development of urban transport systems in many different countries \cite{petersen_behr_2024} \cite{sun2021mobile}. Mobile phones, smartphones in particular have become increasingly accessible, offering options for users across different income levels \cite{worldbank2024P1737510}. Brands such as Tecno, Infinix, Itel, and Samsung’s budget-friendly A-series provide affordable devices for a broad population, while high-end smartphones from Apple, Samsung, and Google Pixel cater to users seeking premium technology \cite{mike2025tecno_infinix_itel} \cite{hogarty2025cheapphones} \cite{lanlehin2025_tecno_infinix_cheap}. In addition to smartphones, wearable devices such as smartwatches also support the development of transport systems \cite{BARKA2024101098}. They also offer both affordable and high-end options. These mobile devices offer a platform for lightweight ride hailing applications that can communicate real time information on vehicle location and availability \cite{preciado‑ortiz2021mobile_apps_transport}. They provide GPS support that allows vehicle tracking and route optimization \cite{SoftwareAdviceGPSVehicleTracking2025}. They also accommodate internet connectivity to allow users to access various ride hailing services \cite{cao2022built_environment}.

In addition, cashless payment systems enabled by technologies like Near Field Communication (NFC) and Radio-Frequency Identification (RFID, allow commuters to pay fares using their debit or credit cards, mobile wallets or even smart cards \cite{ivanova2021_digitization_payments_public_transport}.  Examples include contactless cards like Visa and Mastercard, mobile payment apps like Apple Pay, Google Pay, MTN Mobile Money and Airtel Money and on-board sales units that accept card payments. These cashless systems reduce transaction times, improve convenience, and minimise the need to handle physical money \cite{lexingham2024_cashless_travelling}. Modern ride-hailing applications typically offer multiple payment options, allowing users to select their preferred method \cite{apta2016shared_mobility}. For instance, Uber supports payment via credit/debit cards, and in Uganda users can also pay in cash by selecting the “cash” option in the app \cite{uber2025_payment_options_uganda}. SafeBoda uses its in app Wallet, which can be topped up via mobile money (e.g. MTN, Airtel) or credit or debit cards \cite{kanamanzi2019_safeboda_topup}. Bolt allows in app payments via credit or debit cards or mobile wallets\cite{bolt2025_payments}. 

Improvements in mobile internet connectivity have significantly enhanced global access to ride-hailing and transport-digitization platforms \cite{factmr_ride_hailing_2023}. The expansion of 3G and 4G networks and the gradual introduction of 5G by telecom providers such as MTN, AT and T, Verizon, and Airtel has increased mobile internet speeds and network reliability \cite{IEEE}. High speed connectivity is essential for the operation of services like Uber and Lyft, enabling instant drivers and passengers to be matched instantly, real-time GPS-based vehicle tracking, and smooth processing of digital payments \cite{apta2016shared_mobility}. Navigation platforms such as Google Maps and Waze rely on continuous mobile data connections to deliver live traffic updates, suggest alternative routes, and provide precise turn-by-turn navigation, ultimately reducing travel times and fuel consumption for drivers \cite{stopher2025_waze_revolutionizing_navigation} \cite{kan2018estimating}.

Consistent mobile internet access plays an important role in supporting modern transport information systems by enabling users to access real-time travel information such as schedules, service updates, and digital ticketing options \cite{goto2002new_passenger_support_system} \cite{uitp2022_ticketing_maas}. In addition, improved mobile connectivity has enabled transport and logistics operators to adopt systems that support remote monitoring of vehicle location, speed, and operational status.
Commercial transport operators increasingly use such technologies to enhance operational efficiency, improve service coordination, and support better decision-making. These capabilities also contribute to improved accountability and more structured transport management practices within transport and logistics networks \cite{Ghosh_2025} \cite{KURNIADI}.

Cloud-based systems are a key enabler in the digitization of the global transport sector, providing scalable and flexible infrastructure for managing real-time data, optimizing operations, and improving user experience \cite{Modeshift_CloudTransit_2024}. In modern transport and logistics systems, cloud computing supports the processing of large volumes of live data, including vehicle location through GPS, traffic conditions, booking activity, and system interactions \cite{Ridewyze_RideHailingCloud}.
This capability enables real-time vehicle monitoring, improved coordination between users and operators, and more accurate travel information such as estimated arrival times and service availability. In addition, cloud-based architectures support efficient data sharing between different system components, making them suitable for transport and logistics applications that require continuous updates and multi-user access.

Demand for digital transport services often fluctuates depending on peak travel periods, holidays, and special events. Cloud-based platforms address this variability by providing automatic scalability, allowing computing resources to expand or reduce based on system demand \cite{Lucidchart_ReliabilityAvailability} \cite{Appicial_TaxiZambia_2024}. This ensures consistent system performance, prevents overload, and reduces the need for extensive physical infrastructure.

In addition, cloud-enabled transport systems can support intelligent data processing capabilities such as route optimization by utilizing real-time traffic information and other contextual data. This helps improve travel efficiency by reducing delays and enhancing operational planning \cite{Cloudester_AIRouteOptimization}. Furthermore, automated allocation mechanisms within such systems can streamline operations by matching available transport resources with incoming service requests, improving response time and overall coordination \cite{Ridewyze_AutomatedScheduling}.

Cloud-based analytics platforms enable real-time analysis of demand and supply data, allowing transport systems to respond dynamically to changing user requirements. This capability supports operational adjustments that improve service allocation efficiency and enhance overall system performance \cite{Waehner_DynamicPricing_Kafka}. In transport and logistics contexts, such data-driven approaches can help improve resource distribution by identifying areas of high demand and optimizing the utilization of available vehicles.

In addition, cloud environments support continuous integration and deployment practices, enabling rapid feature updates and system improvements without significant downtime \cite{Cloudfresh_CICD_DevOps}. This ensures that transport applications remain adaptable, maintainable, and capable of evolving in response to user needs and operational challenges.

Cloud-based platformssupport vehicle health tracking, contributing to better maintenance planning and regulatory compliance \cite{Moiboo_VehicleDiagnostics_USA}. Additionally, cloud providers offer robust security features such as encryption, intrusion detection, and multi-factor authentication to protect sensitive user and payment data \cite{Amola}. They also assist transport platforms in complying with international data protection standards, such as the General Data Protection Regulation (GDPR) \cite{GoogleCloudGDPR2025}.

\section{Existing Private Ride-hailing Platforms}
Having established the technological foundations that enable digitized transport solutions, it is instructive to examine how these technologies have been applied in practice through local and international  platforms operating in Uganda. For these digital platforms to succeed in Uganda, they must navigate a complex interplay of opportunities and challenges that differ markedly from more developed markets. Below are some explorations into existing research and case studies on transport and courier apps operating within Uganda, shedding light on both their transformative potential and the persistent hurdles they face. Through this examination, the fundamentals that underpin effective app-based public transport services in Uganda become clearer, guiding the development objectives for Traxis. 

\subsection{Uber}
Uber technologies, is a global transportation company with its headquarters in San Francisco, California, USA. It
develops, markets and operates the Uber car transportation and food delivery mobile applications \cite{uber2026about}. Uber drivers use their own cars or they can rent a car to drive with Uber after passing a criminal background check by Interpol \cite{uber2026rentalfaq}. The Uber application software requires drivers to have smartphones and users must have access to either a smartphone or a mobile website. A user provides his or her destination and then requests a ride; Uber performs a geo-location
search for the nearest driver available and lets a user confirm the request indicating the fare to be paid. When a driver accepts the request, a notification is received and a user can then monitor the car movement on the maps
fragment of the application. Details like name of driver, vehicle registration, make of vehicle and approximated arrival time is provided to the user \cite{uber2026howitworks}.

Uber’s introduction to Kampala marked a significant innovation in the urban transport sector by formalizing the taxi experience through smartphone technology \cite{Kayi2019}. Before Uber, commuting largely depended on informal taxis with limited regulation and inconsistent service quality. Uber’s model introduced on-demand rides, GPS-based tracking, and cashless payments, dramatically enhancing convenience and transparency\cite{Kayi2019}. 

However, the integration of Uber into Uganda’s transport ecosystem faced several challenges, including regulatory resistance and limited adoption beyond higher-income urban users due to affordability constraints and smartphone accessibility \cite{Kayi2019} \cite{IsaacXXXX}.

Uber’s operational model differs significantly from the requirements addressed by Traxis. While it facilitates point-to-point transport for individual users \cite{Kayi2019}, Traxis is designed to support route-based travel where multiple passengers share a vehicle along predefined intercity routes. In such systems, passengers may board and disembark at different points during the journey, resulting in continuously changing seat availability. Unlike Uber, Traxis is designed to support mid-route passenger coordination, where a seat freed at one point along the journey can be reassigned to another passenger further along the route. This enables better seat utilization and reduces revenue loss for operators.

In addition, pricing models differ significantly. Uber fares are calculated dynamically based on distance and time, which can make long-distance travel costly  \cite{Kayi2019} \cite{IsaacXXXX}, \cite{Monitor2016UberKampala}. In contrast, Traxis supports fixed and more predictable pricing structures for specific routes, making it more suitable and accessible for intercity travel.

\subsection{Safeboda}
SafeBoda is a motorcycle ride-hailing company headquartered in Kampala, Uganda, operating across several cities in East and West Africa. It develops, markets, and operates a mobile application that connects passengers to trained motorcycle taxi (boda boda) riders for urban transportation and delivery services \cite{musinguzi2022safeboda}.

SafeBoda riders use their own motorcycles and must undergo a rigorous onboarding process that includes background checks, road safety training, and customer care training before being allowed onto the platform. Riders are also provided with helmets and reflective jackets to enhance safety for both themselves and passengers \cite{muhindo2025rideculture}.

The SafeBoda application requires riders to have smartphones, while users must have access to either a smartphone or, in some cases, USSD/mobile-based alternatives. A user inputs their pickup location and destination, then requests a ride through the app. The system performs a geo-location search to identify the nearest available rider and provides an estimated fare before the ride is confirmed . For courier services, users can request a rider to pick up and deliver parcels to a specified destination. The system performs a geo-location search to identify the nearest available rider and provides an estimated fare before the request is confirmed \cite{safeboda2026faqs}.

Once a rider accepts the request, the user receives a notification and can track the rider’s movement in real time through the map interface. The application displays key details such as the rider’s name, profile photo, motorcycle registration number, and estimated time of arrival \cite{safeboda2024acceptingrides}.

After the trip, both the rider and the passenger can rate each other on a scale of 1 to 5 based on factors such as safety, professionalism, and overall experience. This rating system helps maintain service quality and accountability.

In Uganda, SafeBoda supports both cash and cashless payments. Users can pay using mobile money services integrated within the app or opt to pay cash at the end of the trip, making the service accessible to a wider population \cite{alqahami2022safeboda}.

While SafeBoda is generally more affordable than normal boda bodas \cite{DignitedUberTaxifySafeBoda}, its pricing can still be costly for the average commuter, particularly over long distances where per-kilometre charges accumulate \cite{ShellFoundationSafeBodaImpact}. While this may be suitable for short urban trips, it becomes increasingly expensive and less predictable over longer distances. For intercity travel, such as journeys from Kampala to Kabale, fare estimates provided at the beginning of a trip may differ from the final cost due to factors such as traffic conditions, route changes, or delays. This reduces cost predictability for users and limits affordability for long-distance travel. In contrast, Traxis is designed to support fixed, route-based pricing models commonly used by private shuttle operators. These pricing structures are predetermined for specific routes, making them more affordable and predictable for passengers.

Moreover, SafeBoda operate on a point-to-point transport model, where a vehicle is assigned to a single trip for a specific passenger \cite{ObwotSafeCar2022}. This model does not support shared, route-based transport where multiple passengers can board and disembark at different points along a journey. In contrast, Traxis is designed to support a route-based shuttle system in which passengers may disembark at intermediate locations, and their seats can be reassigned to other passengers further along the route. This dynamic seat management approach improves vehicle utilization and reduces revenue loss for operators while maintaining affordability for passengers.

\subsection{Taxify}
Taxify is an international transportation network company, founded and headquartered in Tallinn, Estonia. The application enables people to request a ride-hailing from
their smartphones. Once the rider has conrmed the pickup location, requested a ride and the driver has accepted. The rider can see the drivers details and both the driver
and rider rate each other after the ride. Taxify drivers undergo a criminal background check and in person training \cite{taxify2018centralasia}.
Taxify (now bolt) leverages Uganda’s mobile money ecosystems like M-Pesa to facilitate digital payments and attract drivers by offering lower commissions than competitors\cite{aPurple_BoltBusinessModel} \cite{Quartz_TaxifyToBolt}. Its market entry disrupted existing ride-hailing dynamics, providing an affordable alternative and extending services to less affluent users. The transition to the Bolt brand signaled a broader vision that includes multi-modal mobility options including e-scooters, reflecting shifting urban transport preferences. Bolt's strategic emphasis on safety features, local partnerships, and customer education has enhanced user and driver retention. This case highlights the critical role of financial inclusion, regulatory alignment, and diversified service offerings for ride-hailing success in Uganda’s evolving urban transport landscape\cite{Pctech_TaxifyUganda2017}. 

 Bolt operates primarily as a private ride-hailing platform where a single booking typically reserves the entire vehicle for one user and their companions. While the platform allows features such as upto three stops within a trip and account-based ride management, it does not support independent seat-based booking by multiple strangers within the same vehicle \cite{Salmon2023RideHailingMonitor}.

As a result, any form of shared travel on Bolt is coordinated through a single booking entity rather than through a structured system where individual passengers independently reserve seats within a shared vehicle. This means that the ride is pre-organized under one request, and additional passengers are included only through the original booking structure rather than through an open seat allocation system.

In contrast, private shuttle systems operate on a seat-based model where multiple unrelated passengers independently book seats within a shared vehicle operating on a fixed route. This creates a structured coordination system that is fundamentally different from Bolt’s booking model.

Traxis is designed to support this shuttle-based structure by enabling independent seat booking, real-time seat availability updates, and coordinated passenger allocation within a single scheduled journey.





\subsection{Lyft}
Lyft, used in the USA and Canada, is a ride-hailing app like Uber \cite{relawding2021lyft}. Lyft offers a Shared ride option designed to facilitate carpooling by matching passengers traveling along similar routes. In this service, each booking is typically limited to one or two passengers per request. This approach enables cost savings through shared travel \cite{yahoo2022lyftshared}. In contrast, a standard Lyft ride provides a private trip for a single party of up to four passengers, ensuring exclusive use of the vehicle. For larger groups, Lyft XL is available, accommodating up to six passengers in a single booking \cite{lyft2026riderpolicies}. These service tiers allow Lyft to cater to different user needs, balancing affordability, convenience, and capacity.
Lyft also provides courier and delivery services through its “Lyft Delivery” offering, which enables drivers to transport a range of items including groceries, retail goods, auto parts, and medical supplies \cite{pymnts2020lyftdelivery}, \cite{rideshareguy2023lyftdelivery}, \cite{fastcompany2020lyft}. Unlike fully consumer-facing delivery platforms, this service primarily operates through partnerships with businesses rather than direct user requests within the standard Lyft application. Deliveries are typically initiated via integrated platforms or third-party partners, such as Olo, allowing restaurants and other organizations to request on-demand, same-day delivery service \cite{littman2021lyftolo}s. Lyft Delivery is therefore largely structured around business-to-business (B2B) and business-to-consumer (B2C) models, rather than person-to-person (P2P) courier services \cite{wiklund2021lyftstrategy}. The service is available in select cities across the United States and parts of Canada, and focuses mainly on essential deliveries such as food, retail items, and pharmacy products.

Despite its diversification, Lyft’s operational model remains centered on urban ride-hailing and on-demand delivery in highly structured digital environments. It does not support fixed-route, intercity shuttle transport systems where passengers independently book seats within a scheduled vehicle operating along a predefined route.

In contrast, Traxis is designed specifically for private shuttle operations in contexts such as Kampala–Kabale routes, where transport is organized around shared vehicles, fixed departure schedules, and seat-based allocation for independent passengers. Unlike Lyft’s ride-matching or ride-sharing mechanisms, Traxis enables structured seat booking within a single scheduled journey, allowing coordinated management of occupancy for intercity travel.

Furthermore, while Lyft Delivery is largely platform- and partner-driven, Traxis integrates passenger transport and courier services within a single shuttle-based system, enabling direct coordination between operators and users for both seats and parcel movement within the same operational framework.



\subsection{Moovit}
Moovit is primarily a comprehensive public transit navigation and Mobility as a Service (MaaS) app, not just a ride-hailing app like Uber. It acts as a trip planner that combines public transit (buses, trains, subways), bikes, scooters, and Moovit also integrates with ride-hailing services. It aggregates data from multiple transport providers, including buses, trains, and informal transit systems, to assist users in planning efficient journeys. The application utilizes real-time data, GPS tracking, and user-generated updates to provide accurate arrival times, route options, and service alerts. Users input their origin and destination, and the system generates optimal routes by combining different modes of transport, including walking segments where necessary \cite{moovit2026about}.

Moovit also supports features such as live directions, step-by-step navigation, and service disruption notifications, enhancing the overall commuting experience \cite{moovit2026navigating}. In addition, the platform contributes to urban mobility planning by providing data analytics tools to cities and transport authorities \cite{moovit2026about}. 

However, Moovit functions primarily as an information and navigation system rather than an operational transport coordination platform. It assists users in selecting routes but does not directly manage vehicle occupancy, seat allocation, or booking within shared transport vehicles.

In contrast, Traxis is designed as an operational coordination system for private shuttle services operating on fixed intercity routes such as Kampala–Kabale. Instead of only providing route information, Traxis enables direct seat booking, real-time seat availability updates, and passenger allocation within specific shuttle vehicles. It also supports operator-level coordination, including monitoring of vehicle status and passenger demand. Furthermore, while Moovit focuses on aggregating transport data across multiple providers.

\subsection{Transit App}
The Transit App is a mobility application designed to help users plan, navigate, and manage public transportation journeys in real time. The application enables users to enter a destination and view available transport options, including buses, trains, and other nearby mobility services, along with estimated departure times and route alternatives. Once a trip is selected, the app provides step-by-step navigation guidance, helping users move through each stage of their journey from origin to destination \cite{TransitApp2025}\cite{transit2026howtouse}.

A key feature of the Transit App is its real-time trip planning capability, which allows users to compare multiple route options based on travel time, convenience, and service availability. The app also supports live directions through its “GO” feature, which provides continuous guidance during travel, including notifications for when to leave, when to transfer, and when to alight from a vehicle. This ensures that users receive timely alerts throughout their journey without needing to constantly monitor the app \cite{transit2026howtouse}.

Additionally, the Transit App provides service alert notifications, which inform users about disruptions such as delays, route changes, or cancellations. Users can also access nearby departures from their current location, improving convenience when planning spontaneous trips. In some cities, the app integrates additional mobility services such as ride-hailing, bikeshare, and e-scooters, offering a more comprehensive multimodal transport solution \cite{transit2026howtouse}.

However, the Transit App functions primarily as an informational and navigation tool rather than a transport coordination system. It does not manage vehicle occupancy, seat allocation, or booking within specific transport services.

In contrast, Traxis is designed as an operational coordination system for private shuttle services operating on fixed intercity routes such as Kampala–Kabale. Instead of only providing route guidance, Traxis enables users to directly book seats, view real-time seat availability, and coordinate travel within a specific shuttle vehicle.

Furthermore, while the Transit App aggregates information across multiple transport providers to assist in journey planning, Traxis focuses on managing a specific transport operator’s services, integrating seat-based passenger booking, real-time vehicle tracking, and courier coordination within a single structured system.

O
\subsection{SWVL}
Swvl is a technology-driven mass transit platform that aims to transform urban transportation by providing efficient, affordable, and accessible shared mobility solutions. Operating across multiple cities in the Middle East, Africa, Asia, and Latin America, Swvl leverages data-driven systems to improve public transport efficiency and reduce reliance on private vehicles. The platform addresses key urban mobility challenges such as traffic congestion, inefficient routing, limited transport accessibility, and safety concerns by introducing optimized, demand-responsive transit solutions \cite{SWVL2024HolidayInnSiemensBosch}.

Swvl operates through a digital ecosystem that connects passengers to shared buses and vehicles via mobile applications. Using real-time data, adaptive routing algorithms, and predictive demand models, the system designs efficient travel routes, improves vehicle occupancy rates, and reduces travel time. It also introduces virtual stops and flexible routing to enhance first- and last-mile connectivity, making public transport more accessible to users in both densely populated and underserved areas. Additionally, Swvl supports use cases such as corporate employee transportation, university transit, and city-wide public mobility systems \cite{narain2025swvlmobility}.

Beyond improving efficiency, Swvl emphasizes sustainability by reducing traffic congestion and lowering carbon emissions through shared rides and optimized fleet utilization. The platform integrates tools for fleet operators, drivers, and administrators, enabling real-time monitoring, scheduling, and service optimization. Overall, Swvl functions as a mobility-as-a-service (MaaS) platform designed to modernize traditional public transport systems and create safer, more reliable, and more sustainable urban mobility networks.

However, Swvl is primarily designed for structured mass transit systems and operates at a system-wide optimization level, rather than focusing on individual private shuttle operators managing specific intercity routes. Its model assumes centralized fleet management and algorithmic route planning across large-scale transport networks.

In contrast, Traxis is designed for privately operated shuttle services such as Kampala–Kabale routes, where transport is organized around fixed schedules, seat-based booking, and operator-controlled vehicles. Instead of dynamically redesigning routes, Traxis focuses on coordinating passenger bookings, managing seat occupancy in real time, and improving operational visibility for individual shuttle operators.

Furthermore, while Swvl optimizes transport at a network level using predictive algorithms, Traxis operates at the service-provider level, enabling direct coordination between passengers and a specific shuttle operator, including features such as seat allocation, real-time vehicle tracking, and integrated courier coordination.

\subsection{Via}
Via is a mobility-as-a-service (MaaS) platform that provides on-demand shared transportation solutions designed to improve urban transit efficiency \cite{ViaIreland2025}. The service allows users to book rides through a mobile application by entering their pickup location and destination. Once a request is made, the system matches passengers traveling in similar directions and assigns them to shared vehicles, reducing travel costs and improving vehicle occupancy \cite{via2026howitworks}.

A key feature of the Via App is the use of designated “virtual bus stops,” where passengers are directed to nearby pickup points instead of exact door-to-door collection \cite{via2026howitworks}. This approach minimizes route deviations, reduces travel time, and increases overall system efficiency. The platform also provides real-time tracking, estimated arrival times, and live trip updates to enhance user convenience and transparency.

Via’s operations are powered by a proprietary algorithm known as “ViaAlgo,” which analyzes traffic conditions, passenger demand, and vehicle locations to optimize routing and minimize waiting times. Unlike traditional ride-hailing services, Via primarily collaborates with public transport agencies and municipalities to complement existing transit systems, offering a flexible middle-ground between fixed-route buses and private taxis. Drivers receive turn-by-turn navigation and optimized pickup instructions through a dedicated driver application \cite{via2026howitworks}.

owever, Via operates primarily as an algorithm-driven shared mobility system designed for dynamic ride pooling and demand-responsive transit. It does not focus on operator-managed private shuttle services where fixed routes, scheduled departures, and structured seat allocation are central to operations.

In contrast, Traxis is designed for privately operated shuttle services such as Kampala–Kabale routes, where transport is organized around predefined schedules, seat-based booking, and operator-controlled vehicles. Rather than dynamically generating ride pools, Traxis provides structured coordination of passengers within a specific shuttle, ensuring clear seat allocation and operational control for the service provider.

Furthermore, while Via emphasizes algorithmic matching and virtual stop optimization at a system level, Traxis focuses on operator-level management, including real-time seat occupancy tracking, passenger booking per seat, vehicle monitoring, and integrated courier coordination within the same shuttle operation.

\subsection{Grab}
Grab is a multi-service mobility and digital platform operating across several Southeast Asian countries, providing ride-hailing, food delivery, and financial services through a unified mobile application. The platform connects users with nearby drivers by using real-time geo-location technology to identify the closest available driver. Once a ride is requested, the application displays essential driver and vehicle information, including the driver’s photo, vehicle details, and license plate number, to ensure safety and easy identification \cite{Wu2025GrabBusinessModel}.

The Grab application offers flexibility by allowing users to either request immediate pickups or schedule rides in advance depending on their travel needs. This enhances convenience and supports both planned and on-demand transportation requirements. In addition, the platform includes an in-app messaging feature that enables communication between drivers and passengers, with built-in translation support to facilitate interaction across different languages \cite{jugnoo2022grabmodel}.

However, Grab operates primarily as a super-app focused on individual ride-hailing and on-demand service delivery across multiple sectors. While it integrates several services within one platform, its transport model remains centered on point-to-point trips between a passenger and a driver, rather than structured coordination of shared, route-based transport systems.

In contrast, Traxis is designed specifically for private shuttle operations such as Kampala–Kabale routes, where multiple passengers share a scheduled vehicle operating along a fixed intercity route. Instead of focusing on individual ride matching, Traxis enables structured seat-based booking, real-time seat occupancy tracking, and coordinated passenger allocation within a single transport service.

Furthermore, while Grab integrates multiple consumer services within a super-app ecosystem, Traxis focuses narrowly on optimizing shuttle operations by combining passenger transport and courier coordination within a single structured system designed for operator-managed, route-based mobility.


\subsection{Comparative analysis table of the different ride-hailing platforms}
The table below provides a summary of the features of the different platforms discussed in the sections above.
% \newgeometry{left=0.5cm, right=0.5cm}
% \begin{longtable}{|p{1.5cm}|p{2cm}|p{2cm}|p{2cm}|p{2cm}|p{2cm}|p{2cm}|p{2cm}|}
% \caption{Comparative Analysis of Transportation Platform Features} \\
% \hline
% Feature & Target market & Courier intergration & Real-time seat availability & Waiting uncertainty & Real-time passenger location & Payment model & Price point \\
% \hline
% \endfirsthead
% \caption{Comparative Analysis of Transportation Platform Features (continued)} \\
% \hline
% Feature & Target market & Stage-based matching & Real-time seat availability & Waiting uncertainty & Real-time passenger location & Payment model & Price point \\
% \hline
% \endhead
% \hline
% \endfoot
% \hline
% \endlastfoot
% Traxis & 15-seater taxis & Real-time stage queue visibility & Provides real-time seat availability information to passengers & Passengers see distance to vehicle via real-time tracking using GPS and maps & Drivers and conductors view real-time passenger hotspots to move towards higher demand & Both cash and cashless payment methods are accepted & Generally affordable for average users \\
% \hline
% Uber & Private cars and motorcycles & Strictly point-to-point & No seat availability info & Passengers see distance to vehicle via real-time tracking using GPS and maps & It follows a passenger initiated model where the passenger inputs their destination and location and wait on a driver to accept their request & Both cash and cashless payment methods are accepted & Relatively expensive, users may share fares \\
% \hline
% Safeboda & Started as motorcycle platform, now implements the use of private cars too & Strictly point-to-point & No seat availability info & Passengers see distance to vehicle via real-time tracking using GPS and maps & It follows a passenger initiated model where the passenger inputs their destination and location and wait on a driver to accept their request & Both cash and cashless payment methods are accepted & Relatively expensive, users may share fares \\
% \hline
% Taxify & Private cars, motorcycles, e-scooters & Strictly point-to-point & No seat availability info & Passengers see distance to vehicle via real-time tracking using GPS and maps & It follows a passenger initiated model where the passenger inputs their destination and location and wait on a driver to accept their request & Both cash and cashless payment methods are accepted & Relatively expensive, users may share fares \\
% \hline
% Lyft & Private cars but has an option where different users can share a ride & Strictly point-to-point & No seat availability info & Passengers see distance to vehicle via real-time tracking using GPS and maps & It follows a passenger initiated model where the passenger inputs their destination and location and wait on a driver to accept their request & Cashless only & Relatively expensive, users may share fares \\
% \hline
% Moovit & Buses and trains & Predefined stops for public transport & Shows general crowdedness, not exact seats & Tracks proximity & Not passenger-initiated, has on-demand features & Cashless preferred, cash for some systems & Affordable, optional premium features \\
% \hline
% Transit app & Buses, trains, e-scooters, bicycles & Predefined stops & Uses capacity predictions, crowdedness & Tracks proximity & Not passenger-initiated, just finds slots & Cashless only & Highly affordable \\
% \hline
% SWVL & Buses and shuttles & Predefined stops & Shows trips with available capacity & Tracks proximity & Passenger requests confirmed by system based on seat availability & Cashless only & Highly affordable \\
% \hline
% Via & Buses & Stops based on user-input destinations & Guarantees seating via booking & Tracks proximity & Passenger-initiated, AI-designated stops & Cashless only & Highly affordable \\
% \hline
% Grab & Private cars and motorcycles & Strictly point-to-point & No seat availability info & Passengers see distance to vehicle & It follows a passenger initiated model where the passenger inputs their destination and location and wait on a driver to accept their request & Cashless only & Highly affordable \\
% \hline
% \end{longtable}
% \restoregeometry

\begingroup
\begin{longtblr}[
  caption = {Comparative Analysis of Transportation Platform Features},
  label = {tab:transport-comparison}
]{
  colspec = {|p{1.5cm}|p{1.5cm}|p{2cm}|p{2cm}|p{2cm}|p{2cm}|p{1.8cm}|p{1.7cm}|},
  rowhead = 1
}
\hline
Feature &
Transport Model &
Booking structure &
Route structure &
Seat management  &
Operational coordination &
Pricing model  &
Courier integration \\
\hline
Traxis &
Fixed-route private shuttle coordination &
Independent seat-based booking &
Fixed intercity routes (e.g., Kampala–Kabale) &
Real-time seat allocation and  tracking &
Operator-level shuttle management &
Fixed fare per seat and  courier pricing &
Fully integrated passenger and  parcel system  \\
\hline
Uber &
Point-to-point rides or navigation support &
Whole ride per request or trip planning only &
Dynamic or user-defined &
PNot exposed to users &
Driver-passenger matching or route guidance &
Dynamic or distance-based &
Separate or limited \\
\hline
Safeboda &
Point-to-point rides or navigation support &
Whole ride per request or trip planning only &
Dynamic or user-defined &
PNot exposed to users &
Driver-passenger matching or route guidance &
Dynamic or distance-based &
Separate or limited \\
\hline
Taxify (Bolt) &
Point-to-point rides or navigation support &
Whole ride per request or trip planning only &
Dynamic or user-defined &
PNot exposed to users &
Driver-passenger matching or route guidance &
Dynamic or distance-based &
Separate or limited \\
\hline
Lyft &
Point-to-point rides or navigation support &
Whole ride per request or trip planning only &
Dynamic or user-defined &
PNot exposed to users &
Driver-passenger matching or route guidance &
Dynamic or distance-based &
Separate or limited \\
\hline
Moovit &
Point-to-point rides or navigation support &
Whole ride per request or trip planning only &
Dynamic or user-defined &
PNot exposed to users &
Driver-passenger matching or route guidance &
Dynamic or distance-based &
Separate or limited \\
\hline
Transit App &
Point-to-point rides or navigation support &
Whole ride per request or trip planning only &
Dynamic or user-defined &
PNot exposed to users &
Driver-passenger matching or route guidance &
Dynamic or distance-based &
Separate or limited \\
\hline
SWVL &
Shared mass transit (bus-like systems) &
Grouped/shared ride matching &
Algorithm-optimized or flexible routes &
Limited grouping, not seat-by-seat control &
System-level fleet optimization &
Variable/shared fares &
Not core functionality \\
\hline
Via &
Shared mass transit (bus-like systems) &
Grouped/shared ride matching &
Algorithm-optimized or flexible routes &
Limited grouping, not seat-by-seat control &
System-level fleet optimization &
Variable/shared fares &
Not core functionality \\
\hline
Grab &
Point-to-point rides or navigation support &
Whole ride per request or trip planning only &
Dynamic or user-defined &
PNot exposed to users &
Driver-passenger matching or route guidance &
Dynamic or distance-based &
Separate or limited \\
\hline
\end{longtblr}

\endgroup

\section{Conclusion}
In summary, the growth of ride-hailing and digital mobility platforms has marked a significant transformation in urban transportation systems globally and within Uganda. International platforms such as Uber have introduced digital convenience through features such as GPS tracking, cashless payments, and formalized booking systems, helping to modernize urban taxi experiences. However, their adoption in Uganda remains largely concentrated among higher-income urban users, limiting their inclusivity within the broader transport population.

Locally adapted platforms such as SafeBoda have demonstrated the importance of contextual design by integrating mobile payments, safety training, and affordability into their service model, ensuring wider usability within Uganda’s transport ecosystem. Similarly, platforms like Bolt have expanded access through competitive pricing and mobile money integration, further contributing to the evolution of app-based mobility services in the region.

Despite these advancements, existing platforms primarily focus on either private point-to-point transport or generalized ride-hailing services, with limited support for structured coordination of shared, fixed-route shuttle systems that dominate intercity travel in Uganda. This creates a gap in how seat-based occupancy, scheduled departures, and operator-managed shuttle services are coordinated digitally.

These observations highlight the need for a context-aware and system-specific solution tailored to Uganda’s private shuttle transport sector. Traxis is therefore proposed as a digital platform designed to improve coordination, communication, and operational efficiency within fixed-route shuttle services through real-time seat management, structured booking, and integrated transport–courier functionality.